\section{Carga transportada y perturbaciones externas}
\label{sec:perturbaciones}

\subsection{Carga transportada}
La masa del sólido rígido pasa a ser la masa total
\begin{equation}
  m_{total} = m_d + m_c ,
\end{equation}
donde $m_c$ es el parámetro \texttt{m\_payload}, igual a $0$ por defecto. La carga se
supone concentrada en el centro de gravedad, por lo que no modifica el tensor de
inercia $J$. El bloque \texttt{ForceCompute} del controlador sigue usando la masa
nominal $m_d$, es decir, el controlador \emph{desconoce} la carga. De este modo se
puede estudiar la variación $m \rightarrow m + \Delta m$ (Modelado, ec.~41) y
comprobar que la acción integral del lazo de velocidad compensa el peso añadido.

La traslación del cuerpo queda, en el sistema inercial,
\begin{equation}
  m_{total}\, \dot{\mathbf{v}} = \mathbf{F}_T + \mathbf{F}_g + \mathbf{F}_d + \mathbf{F}_{wind},
  \qquad
  \mathbf{F}_g = \begin{bmatrix} 0 & 0 & -m_{total}\, g \end{bmatrix}^T .
\end{equation}

\paragraph{Carga máxima (pregunta 2.1 del Informe~1).}
El empuje máximo está limitado por $\omega_{max}$ y por la tensión de la batería. La
velocidad máxima alcanzable con $V_{bat}$ se obtiene de la condición de régimen
permanente
\begin{equation}
  \frac{R\,k_m}{k_t}\,\omega^2 + \left(k_e + \frac{R\,b_m}{k_t}\right)\omega - V_{bat} = 0 ,
\end{equation}
que con los valores por defecto da $\omega \approx \SI{1042}{rad/s} > \omega_{max}$. El
límite efectivo es, por tanto, $\omega_{max}$:
\begin{equation}
  m_{c,max} = \frac{4\,k_f\,\omega_{max}^2}{g} - m_d .
\end{equation}

\subsection{Resistencia aerodinámica}
Se adopta un modelo lineal en la velocidad inercial:
\begin{equation}
  \mathbf{F}_d = -c_d\, \mathbf{v}
\end{equation}
Por defecto $c_d = 0$, que coincide con la simplificación del Informe~1, y puede
activarse dándole un valor positivo.

\subsection{Viento}
Se modela como una fuerza constante en el sistema inercial que comienza en
$t = t_{wind}$ (Modelado, ec.~40):
\begin{equation}
  \mathbf{F}_{wind}(t) =
  \begin{cases}
    \mathbf{0} & t < t_{wind} \\
    \begin{bmatrix} F_x & F_y & F_z \end{bmatrix}^T & t \ge t_{wind}
  \end{cases}
\end{equation}
Con $F_y = F_z = 0$ se recupera el caso $\mathbf{F}_{wind} = [F_x\;0\;0]^T$ del
documento de modelado.

\subsection{Aplicación sobre el cuerpo}
El bloque \texttt{Perturbaciones} se conecta al \texttt{frame\_a} del sólido rígido y
aplica $\mathbf{F} = \mathbf{F}_d + \mathbf{F}_{wind}$ del mismo modo que
\texttt{MultiBody.Forces.WorldForce}. La fuerza, definida en el sistema inercial, se
proyecta al sistema del conector mediante la matriz de orientación:
\begin{equation}
  \mathbf{f}_{frame} = -R^T\, \mathbf{F}, \qquad \boldsymbol{\tau}_{frame} = \mathbf{0},
\end{equation}
siendo $R$ la matriz de rotación cuerpo $\rightarrow$ inercial. La velocidad se obtiene
como $\mathbf{v} = \dot{\mathbf{r}}_0$ a partir de la posición absoluta del conector.
